Al estudiar distintas ramas de las matemáticas de manera formal es
inevitable notar que existen ciertos patrones que se repiten. Todas
empiezan con los axiomas que caracterizan los objetos de estudio, como
pueden ser los grupos, espacios vectoriales, espacios topológicos,
variedades, etc., y pese a que estos se definen de forma muy
diferente, en la mayoría se pueden definir de manera natural conceptos
como subobjetos, objetos producto u objetos cociente, entre
otros. Además, estos objetos no se suelen estudiar de forma aislada,
sino que se estudian funciones entre los objetos que cumplen ciertas
propiedades, como pueden ser las aplicaciones lineales en el álgebra
lineal, las aplicaciones continuas en la topología y los homomorfismos
del álgebra abstracta.

Estos patrones no son una coincidencia, pues las propiedades generales que
exhiben son las mismas en todos los casos. Por ejemplo, la clase de funciones
<<destacables>> entre objetos es cerrada para la composición, y el producto de
subobjetos de dos objetos es un subobjeto del producto de los objetos
originales. Esto motiva el estudio de dichos patrones como otra área de las
matemáticas, con el fin de obtener propiedades aplicables a muchas de las áreas
de las matemáticas existentes o, incluso, a áreas desarrolladas posteriormente,
como veremos que ocurre con la teoría de la computación. Así, en este área de
las matemáticas los objetos de estudio son representaciones de los conceptos
fundamentales de otras áreas, que podemos representar como la clase de los
objetos de estudio junto con la clase de funciones destacables entre ellos,
llamadas \emph{morfismos}, lo que da lugar al concepto de \emph{categoría}.
Este capítulo introduce las categorías y estudia sus propiedades de manera
general, y se basa principalmente en \cite[caps. 3, 4, 7 y 10]{joyofcats}.

\begin{definition}
  \label{def:category}
  Una \conc{categoría} es una tupla $\cC$ formada por los siguientes elementos:
  \begin{enumerate}
  \item Una clase $\Ob{\cC}$ de \conc{objetos}.
  \item Una clase $\Mor{\cC}$ de \conc{morfismos}.
  \item Dos funciones $\dom,\cod:\Mor{\cC}\to\Ob{\cC}$, llamadas
    respectivamente \conc{dominio} y \conc{codominio}.

    Para $f\in\Mor{\cC}$, escribimos $f:a\to b$ si $\dom{f}=a$ y $\cod{f}=b$, y
    llamamos $\hom_{\cC}(a, b)$ o simplemente $\hom(a, b)$ a la clase de
    morfismos $f:a\to b$. %, que generalmente requeriremos que sea un conjunto.
  \item Una función
    $\circ:\bigcup_{a,b,c}(\hom(b,c)\times\hom(a,b))\to\Mor{\cC}$ que a cada
    $f:a\to b$ y $g:b\to c$ le asigna su \conc{composición} $g\circ f:a\to c$, y
    que debe ser asociativa, es decir, para $f:a\to b$, $g:b\to c$ y $h:c\to d$,
    $h\circ(g\circ f) = (h\circ g)\circ f$.
  \item Una función $1:\Ob{\cC}\to\Mor{\cC}$ que a cada objeto $a$ le asigna la
    \conc{identidad} $1_a:a\to a$, que cumple que, para $f:a\to b$,
    $f=1_b\circ f=f\circ 1_a$.
  \end{enumerate}
\end{definition}

Es común en teoría de categorías representar situaciones con diagramas, como los
usados en buena parte del álgebra, donde los vértices representan objetos y las
flechas representan morfismos, y una flecha aparece punteada si su existencia se
debe a la existencia de las otras flechas del diagrama. No se suelen representar
las flechas identidad ni la composición de flechas que ya aparecen en el
diagrama.

Un diagrama \conc{conmuta} si, dados dos caminos cualesquiera del diagrama con
el mismo objeto de origen y de destino, la composición de los morfismos en cada
camino coincide. Así, por ejemplo, los axiomas de la composición e identidad de
categorías se pueden expresar como en la figura \ref{fig:cat-axiom}.

\begin{figure}
  \centering
  \begin{subfigure}{.45\textwidth}
    \centering
    \begin{diagram}
      \path (0,2) node(A) {$a$} (2,2) node(B) {$b$}
            (0,0) node(C) {$c$} (2,0) node(D) {$d$};
      \draw[->] (A) -- node[above]{$f$} (B);
      \draw[->] (B) -- node[right]{$h\circ g$} (D);
      \draw[->] (A) -- node[left]{$g\circ f$} (C);
      \draw[->] (C) -- node[below]{$h$} (D);
    \end{diagram}
    \caption{Asociatividad}
  \end{subfigure}
  \hfill
  \begin{subfigure}{.45\textwidth}
    \centering
    \begin{diagram}
      \path (0,2) node(A){$a$} (2,2) node(AP){$a$}
            (0,0) node(B){$b$} (2,0) node(BP){$b$};
      \draw[->] (A) -- node[above]{$1_a$} (AP);
      \draw[->] (AP) -- node[right]{$f$} (BP);
      \draw[->] (A) -- node[left]{$f$} (B);
      \draw[->] (B) -- node[below]{$1_b$} (BP);
      \draw[->] (A) -- node[above]{$f$} (BP);
    \end{diagram}
    \caption{Elemento neutro}
  \end{subfigure}
  \caption{Axiomas de categoría}
  \label{fig:cat-axiom}
\end{figure}

\begin{samepage}
  \begin{example}
    Ciertos conceptos fundacionales se pueden ver como categorías:
    \begin{enumerate}
    \item La categoría $\bSet$ tiene como objetos todos los conjuntos y como
      morfismos todas las funciones, cualificadas por su dominio y codominio, con
      la composición e identidad obvias.
    \item Un \conc{preorden} es una relación reflexiva y transitiva, que se puede
      ver intuitivamente como un orden parcial entre clases de equivalencia. Un
      \conc{conjunto preordenado} es un conjunto con un preorden asociado, y
      llamamos $\bPrord$ a la categoría de los conjuntos preordenados cuyos
      morfismos son funciones que conservan el preorden, es decir, funciones
      $f:(a,\preceq)\to(b,\preccurlyeq)$ tales que para $x,y\in a$,
      $x\preceq y\implies f(x)\preccurlyeq f(y)$.
    \item \label{enu:mot-subcat} La categoría $\bOrd$ tiene como objetos los
      conjuntos parcialmente ordenados y como morfismos las funciones que
      conservan el orden. $\bLat$ es similar pero los objetos son retículos y los
      morfismos deben además preservar supremos e ínfimos de pares de elementos.
    \item La categoría $\bGrph$ tiene como objetos los grafos dirigidos y como
      morfismos las funciones entre los vértices de dos grafos que llevan ejes del
      primer grafo a ejes del segundo.
    \end{enumerate}
  \end{example}
\end{samepage}

El apartado \ref{enu:mot-subcat} de la lista anterior lleva naturalmente al
concepto de subcategoría.

\begin{definition}
  Una categoría $\cB$ es una \conc{subcategoría} de una $\cC$ si
  $\Ob{\cB}\subseteq\Ob{\cC}$, $\Mor{\cB}\subseteq\Mor{\cC}$ y las funciones
  dominio, codominio, composición e identidad son restricciones de las
  correspondientes funciones de $\cC$, y es una \conc{subcategoría completa} si
  además, para $a,b\in\Ob{\cB}$, $\hom_\cB(a,b)=\hom_\cC(a,b)$.
\end{definition}

Así, $\bLat$ es una subcategoría no completa de $\bOrd$, mientras que $\bOrd$ es
una subcategoría completa de $\bPrord$, y esta a su vez es una subcategoría
completa de $\bGrph$ si consideramos las relaciones como grafos dirigidos y
permitimos ejes reflexivos.

\section{Categorías algebraicas}

\begin{example}\label{ex:variety}
  En álgebra, muchas categorías se definen de manera similar:
  \begin{enumerate}
  \item $\bSmgrp$, la categoría de los semigrupos con las funciones que
    conservan su operación.
  \item $\bMon$, la subcategoría no completa de $\bSmgrp$ de los monoides con
    las funciones que conservan su operación y elemento identidad.
  \item $\bGrp$, la subcategoría completa de $\bMon$ formada por los grupos y
    sus homomorfismos.
  \item $\bAb$, la subcategoría de $\bGrp$ de grupos abelianos.
  \item $\bRing$, la categoría de anillos unitarios y sus homomorfismos.
    Incluímos aquí el anillo trivial en el que $1=0$.
  \item $\bCRng$, la subcategoría completa de $\bRing$ de los anillos unitarios
    conmutativos.
  \end{enumerate}
\end{example}

En todos estos casos los objetos son conjuntos con una serie de
operaciones, y los morfismos son funciones que respetan esas
operaciones en el sentido evidente. La lista de operaciones se puede
modelar como sigue.\cite[p. 120]{maclane}

\begin{definition}
  Un \conc{conjunto graduado} es un conjunto $I$ de \conc{operadores} junto con
  una función $a:I\to\sNat$ llamada \conc{aridad}.  Una \conc{acción} de $I$ en
  un conjunto $S$ es una familia $\{\mu_i : S^{a(i)} \to S\}_{i\in I}$ de
  operaciones en $S$ asociadas a los operadores de $I$.
\end{definition}

Así, por ejemplo, los operadores en $\bAb$ serían $+$ de aridad 2, $-$
de aridad 1 y $0$ de aridad 0. Queda definir las propiedades de los
operadores.
    
\begin{definition} Sea $(I, a)$ un conjunto graduado:
  \begin{enumerate}
  \item El conjunto de \conc{operadores derivados} de $I$ es el conjunto
    graduado $\Lambda$ formado por las siguientes expresiones:
    \begin{enumerate}
    \item El \conc{operador identidad}, $id$, de aridad 1.
    \item Todos los operadores de $I$ con su aridad en $I$.
    \item Para $\omega\in\Lambda$ de aridad $n$ e $i_1,\dots,i_n\in\Lambda$ de
      aridades $a_1,\dots,a_n$, $\omega(i_1,\dots,i_n)$ de aridad
      $a_1+\dots+a_n$.
    \item Para $\lambda\in \Lambda$ de aridad $n$ y
      $f:\{1,\dots,n\}\to\{1,\dots,m\}$, $\lambda_f$ de aridad $m$.
    \end{enumerate}
  \item Si $\mu$ es una acción sobre $I$, la extensión de $\mu$ a $\Lambda$ es
    la acción $\nu$ sobre $\Lambda$ dada por $\nu_i=\mu_i$ para $i\in I$,
    $nmu_{id}(x)\coloneqq x$,
    $\nu_{\omega(i_1,\dots,i_n)}(x_{11},\dots,x_{1a_1},\dots,x_{n1},\dots,x_{na_n})\coloneqq\nu_\omega(\nu_{i_1}(x_{11},\dots,x_{1a_1}),\dots,\nu_{i_n}(x_{n1},\dots,x_{na_n}))$
    y
    $\nu_{\lambda_f}(x_1,\dots,x_m)\coloneqq\nu_\lambda(x_{f(1)},\dots,x_{f(n)})$.
  \item Una \conc{identidad} sobre $I$ es un par $(\lambda,\sigma)$ de
    operadores $\Lambda$ de igual aridad.
  \item Una acción $\mu$ sobre $I$ \conc{satisface} una identidad
    $(\lambda, \sigma)$ sobre $I$ si $\nu_\lambda=\nu_\sigma$, donde $\nu$ es la
    extensión de $\mu$ a $\Lambda$.
  \end{enumerate}
\end{definition}

Normalmente las identidades $(\lambda, \sigma)$ se representan como igualdades
$\lambda=\sigma$, y $\lambda$ y $\sigma$ se representan de forma obvia como
expresiones que dependen de los operadores base y una serie de parámetros de
entrada, de modo que las propiedades de $\bAb$ son $(x+y)+z=x+(y+z)$, $0+x=x$,
$x+(-x)=0$ y $x+y=y+x$.

\begin{definition}
  Sean $\Omega$ un conjunto graduado finito y $E$ un conjunto finito de
  identidades sobre $\Omega$:
  \begin{enumerate}
  \item Una $(\Omega,E)$-álgebra es un par $(S,\mu)$ formado por un conjunto $S$
    y una acción $\mu$ de $\Omega$ sobre $S$ que satisface las identidades en
    $E$.
  \item Una \conc{variedad algebraica} es una categoría $(\Omega,E)\dash\bAlg$
    cuyos objetos son las $(\Omega,E)$-álgebras y cuyos morfismos
    $(A,\mu)\to (B,\nu)$ son las funciones $f:A\to B$ tales que, para
    $\omega\in\Omega$ de aridad $n$ y $x_1,\dots,x_n\in A$,
    $f(\mu_\omega(x_1,\dots,x_n))=\nu_\omega(f(x_1),\dots,f(x_n))$.
  \end{enumerate}
\end{definition}

Claramente todas las categorías del ejemplo \ref{ex:variety} son variedades
algebraicas, y de hecho podemos ver $\bSet$ como
$(\emptyset,\emptyset)\dash\bAlg$, pero sin embargo $\bField$, la subcategoría
completa de $\bRing$ de los cuerpos (incluyendo el cuerpo trivial), no es una
variedad algebraica, ya que requiere una propiedad de la inversa del producto,
que no está definida en el 0.

Otra categoría interesante es $R\dash\bMod$, la clase de módulos de un anillo
conmutativo $R$ y homomorfismos de $R$-módulos. $\sInt\dash\bMod$ es
esencialmente $\bAb$ y, para un cuerpo $K$, $K\dash\bMod$ es la categoría de
$K$-espacios vectoriales, por lo que la escribimos como $K\dash\bVec$ o
simplemente $\bVec$ si $K=\sReal$.

\section{\label{sec:cat-abstract}Categorías abstractas}

Hasta ahora, en todas las categorías que hemos visto, los objetos son conjuntos
con estructura y los morfismos son funciones entre ellos. Las categorías de esta
forma se llaman \emph{constructos}\footnote{Veremos una definición más abstracta
  de constructo cuando veamos los funtores.}, y aunque son muy comunes, también
hay muchas categorías relevantes que no son constructos.

\begin{example}
  Para un anillo $R$, la categoría $R\dash\bMat$ tiene como objetos los números
  naturales, como morfismos de $n$ a $m$ las matrices de tamaño $m\times n$,
  como composición el producto de matrices y como identidad la matriz identidad
  del tamaño correspondiente.
\end{example}

\begin{example}
  Algunas estructuras matemáticas se pueden ver como categorías.
  \begin{enumerate}
  \item Una categoría es \conc{discreta} si sólo tiene los morfismos
    identidad. Un conjunto, o en general una clase, se puede ver como una
    categoría discreta cuyos objetos son los elementos del conjunto.
  \item Una categoría es \conc{fina} si sus \conc{conjuntos hom} (conjuntos de
    la forma $\hom(a,b)$) tienen como máximo un elemento. Un conjunto
    preordenado $(X,\preceq)$ se puede ver como una categoría fina cuyos objetos
    son los elementos del conjunto y tal que, para $x,y\in X$, $\hom(x,y)$ está
    habitado si y sólo si $x\preceq y$.
  \item Un monoide se puede ver como una categoría con un solo objeto, cuyos
    morfismos son los elementos del monoide, la identidad es su elemento neutro
    y la composición es el producto.
  \end{enumerate}
\end{example}

\begin{example}
  Las siguientes categorías se usan principalmente en el estudio de categorías
  más complicadas:
  \begin{enumerate}
  \item La categoría vacía, $\bZero$, sin objetos.
  \item La categoría discreta unipuntual, $\bOne$.
  \item La categoría discreta de dos objetos, $\bTwo$.
  \item La categoría $\bDown$, con dos objetos y una sola flecha de uno a otro
    ($\bullet\to\bullet$).
  \end{enumerate}
\end{example}


\section{Categorías topológicas y analíticas}

La principal categoría topológica es $\bTop$, formada por los espacios
topológicos y las funciones continuas entre ellos.

\begin{example}
  Algunos constructos usados en topología tienen los mismos objetos pero
  distintas clases de morfismos, permitiendo estudiar los objetos desde
  distintas perspectivas. Por ejemplo, las siguientes tres categorías tienen
  como objetos los espacios métricos:
  \begin{enumerate}
  \item $\bMetc$, con las funciones continuas.
  \item $\bMetu$, con las funciones uniformemente continuas.
  \item $\bMet$, con las contracciones, funciones continuas que
    \emph{acercan} los puntos.
  \end{enumerate}
  Asimismo tenemos la categoría $\bBanb$ de los espacios de Banach con las formas
  lineales acotadas y $\bBan$ con las contracciones lineales.
\end{example}


\begin{example}
  Sea $X$ un espacio topológico.  Recordemos que, para $x,y\in X$, un camino de
  $x$ a $y$ es una función continua $f:[0,1]\to X$ con $f(0)=x$ y $f(1)=y$, y
  que dos caminos $f$ y $g$ de $x$ a $y$ son homotópicamente equivalentes si
  existe $F:[0,1]\times[0,1]\to X$ tal que, para $s,t\in[0,1]$, $F(t,0)=f(t)$,
  $F(t,1)=g(t)$, $F(0,s)=x$ y $F(1,s)=y$. Entonces el \conc{grupoide
    fundamental}\cite[p. 20]{maclane} de $X$ es una categoría $\pi(X)$ cuyos
  objetos son los puntos de $X$ y tal que $\hom(x,y)$ es el conjunto cociente de
  caminos de $x$ a $y$ por equivalencia homotópica, tomando la concatenación de
  (clases de) caminos como composición y la clase del camino constante como
  morfismo identidad.

  Para $x\in X$, $\hom(x,x)$ es un grupo con la composición, y si existe un
  camino $f:x\to y$, $\hom(x,x)$ y $\hom(y,y)$ son isomorfos por el isomorfismo
  $g\mapsto fgf^{-1}$, definiendo $f^{-1}$ de la forma obvia, con lo que si $X$
  es conexo por arcos, $\hom(x,x)$ es el mismo para todo $x$ salvo
  isomorfismo, y se llama \conc{grupo fundamental} de $X$. Esto permite tratar
  espacios topológicos de manera algebraica.
\end{example}

\section{Isomorfismos}

Ya hemos visto las categorías más importantes de distintas áreas, por lo que
pasamos a ver algunas propiedades de sus objetos y morfismos. Un primer ejemplo
es el de isomorfismo, que se suele definir como un homomorfismo biyectivo cuya
inversa también es un homomorfismo. En topología tenemos también el concepto
de homeomorfismo, que viene a ser lo mismo salvo por el cambio de terminología.

\begin{definition}
  Un morfismo $f:a\to b$ es un \conc{isomorfismo}, si existe otro morfismo
  $g:b\to a$ tal que $g\circ f=1_a$ y $f\circ g=1_b$, en cuyo caso escribimos
  $f:a\cong b$ y decimos que $a$ y $b$ son \conc{isomorfos}, $a\cong b$.
\end{definition}

Claramente el inverso de un isomorfismo es único, pues si $f:a\to b$ tiene
inversos $g,h:b\to a$, entonces $g=g\circ 1_b=g\circ f\circ h=1_a\circ h=h$.
Llamamos entonces $f^{-1}$ a la inversa de $f$. Además, la composición de dos
isomorfismos es un isomorfismo.

\begin{example}
  Algunos isomorfismos:
  \begin{enumerate}
  \item En $\bSmgrp$, $\bMon$, $\bGrp$, $\bAb$, $\bRing$ y $R\dash\bMod$, el
    concepto de isomorfismo categórico se corresponde con el concepto usual de
    isomorfismo.
  \item En $\bTop$ y $\bMetc$, los isomorfismos son los homeomorfismos.
  \item En $\bBanb$, los isomorfismos son los isomorfismos topológicos, es
    decir, los isomorfismos vectoriales que son homeomorfismos.
  \item En $\bMet$ los isomorfismos son las isometrías, mientras que en $\bBan$
    son los isomorfismos isométricos.
  \item En $\bSet$, los isomorfismos son las funciones biyectivas.
  \item Un \conc{grupoide} es una categoría en la que todos los morfismos son
    isomorfismos. Esto ocurre, por ejemplo, en el grupoide fundamental.
  \item En un conjunto preordenado visto como categoría, dos objetos $a$ y $b$
    son isomorfos si y sólo si $a\preceq b\preceq a$, por lo que en particular
    un conjunto (parcialmente) ordenado visto como categoría no tiene
    isomorfismos salvo las identidades.
  \item En un monoide visto como categoría, los isomorfismos son los elementos
    invertibles.
  \item En $R\dash\bMat$, los isomorfismos son las matrices invertibles.
  \end{enumerate}
\end{example}

\section{Objetos iniciales y finales}

En álgebra, muchas categorías tienen un objeto trivial, como el cuerpo trivial,
el grupo trivial o el espacio vectorial trivial. Sin embargo, a la hora de
intentar definir la topología trivial nos encontramos que hay dos candidatos
posibles, la topología vacía y la unipuntual, y ninguna es enteramente
satisfactoria.  Esto se debe a que, al hablar de un objeto trivial estamos
juntando dos propiedades: por un lado, el objeto trivial está contenido en todos
los demás, y por otro, siempre es posible pasar de un objeto al objeto trivial
por un morfismo.

\begin{definition}
  Un objeto $a$ es \conc{inicial} si para cualquier otro objeto $x$ existe un
  único morfismo $f_x:a\to x$, es \conc{final} si para cualquier otro objeto $x$
  existe un único morfismo $g_x:x\to a$, y es \conc{cero} si es inicial y final.
\end{definition}

\begin{example}
  Veamos los objetos iniciales y finales de algunas categorías típicas.
  \begin{enumerate}
  \item En $\bMon$, $\bGrp$ y $\bAb$, el grupo trivial es un objeto cero; en
    $\bRing$ y $\bField$ lo es el cuerpo trivial, y en $R\dash\bMod$ lo es el
    módulo trivial.
  \item En $\bSet$, el conjunto vacío $\emptyset$ es inicial, mientras que los
    conjuntos unipuntuales $\{*\}$ son finales. Lo mismo ocurre en $\bTop$ y en
    $\bOrd$ dotando a $\emptyset$ y $\{*\}$ de la única estructura posible en
    cada caso.
  \item En un conjunto ordenado visto como categoría, un objeto inicial es un
    mínimo y un objeto final es un máximo. En particular el único conjunto
    ordenado con un cero es el unipuntual.
  \item En $R\dash\bMat$, el 0 es un cero, pues para cada $m$ hay una única
    matriz de tamaño $0\times m$ y una de tamaño $m\times 0$.
  \end{enumerate}
\end{example}

Podríamos preguntarnos si categorías como $\bSet$ o $\bTop$ tienen un cero.
La respuesta es que no, como se deduce de las siguientes proposiciones.

\begin{proposition}
  Los objetos iniciales son únicos salvo isomorfismo:
  \begin{enumerate}
  \item Si $a$ y $b$ son objetos iniciales (de la misma categoría), entonces son
    isomorfos.
    \begin{proof}
      Por hipótesis existen $f:a\to b$ y $g:b\to a$, pero como sólo hay un morfismo
      $a\to a$, este debe ser $1_a$ y por tanto $g\circ f=1_a$, y análogamente
      $f\circ g=1_b$, con lo que $f$ es un isomorfismo.
    \end{proof}
  \item Si $a\cong b$ y $a$ es inicial, entonces $b$ es inicial.
    \begin{proof}
      Sea $f:b\to a$ un isomorfismo y $x$ un objeto cualquiera, existe un único
      $h:a\to x$ y por tanto $h\circ f$ es un morfismo de $b$ a $x$, que es
      único ya que, para $k:b\to x$, $k\circ f^{-1}:a\to x$ y así
      $k\circ f^{-1}=h$, con lo que $k=h\circ f$.
    \end{proof}
  \end{enumerate}
\end{proposition}

La siguiente proposición se prueba de forma análoga.

\begin{proposition}
  Los objetos finales son únicos salvo isomorfismo:
  \begin{enumerate}
  \item Si $a$ y $b$ son objetos finales, entonces son isomorfos.
  \item Si $a\cong b$ y $a$ es final, entonces $b$ es final.
  \end{enumerate}
\end{proposition}

Así, como $\emptyset$ y $\{*\}$ no son isomorfos en $\bSet$, $\bTop$ ni $\bOrd$
(ni en ningún otro constructo), estas categorías no tienen objetos cero.  En
general diremos que un objeto con una cierta propiedad es \conc{único salvo
  isomorfismo} si los objetos que tienen esa propiedad son precisamente los que
son isomorfos a dicho objeto.

\begin{corollary}
  En los grupoides no vacíos, las siguientes condiciones son equivalentes:
  \begin{enumerate}
  \item \label{enu:goid-initial} Existe un objeto inicial.
  \item \label{enu:goid-final} Existe un objeto final.
  \item \label{enu:goid-zero} Existe un objeto cero.
  \item \label{enu:goid-all} Todos los objetos son cero.
  \end{enumerate}
  En particular, en el grupoide fundamental $\pi(X)$ (con $X$ no vacío), esto
  ocurre si y sólo si $X$ es simplemente conexo.
\end{corollary}
\begin{proof}
  La equivalencia (\ref{enu:goid-initial}$\iff$\ref{enu:goid-final}) se debe a
  que, por unicidad del isomorfismo inverso, $\hom(a,b)=\hom(b,a)$ para
  cualesquiera $a$ y $b$, y esto prueba la equivalencia con
  (\ref{enu:goid-zero}). Ahora bien, si $a$ es cero y $x$ es otro objeto, el
  único morfismo $a\to x$ es un isomorfismo, por lo que $x$ es cero y queda
  probado (\ref{enu:goid-zero}$\iff$\ref{enu:goid-all}).  La última
  observación es por definición.
\end{proof}

\section{Monomorfismos y epimorfismos}

En muchas ramas del álgebra llamamos monomorfismos a los morfismos inyectivos y
epimorfismos a los morfismos suprayectivos. Esta definición depende de que los
morfismos sean funciones, por lo que no nos sirve para categorías generales.

Para los monomorfismos podemos basarnos en que, en $\bSet$, los elementos de un
conjunto $S$ se identifican con los morfismos $\{*\}\to S$, con lo que la
propiedad de que $f(x)=f(y)\implies x=y$ se puede traducir como sigue.

\begin{definition}
  Un morfismo $f:a\to b$ es un \conc{monomorfismo} (escrito $f:a\monicTo b$) si
  es cancelable por la izquierda, es decir, si para cualesquiera $h,k:c\to a$
  con $f\circ h=f\circ k$ se tiene $h=k$.
\end{definition}

\begin{example}[Monomorfismos en constructos]\;
  % En el caso de $\bSet$, tomando $c=\{*\}$ se tiene por construcción que todo
  % monomorfismo es inyectivo, y recíprocamente, si $f:X\to Y$ es una función
  % inyectiva y $h,k:S\to X$ con $f\circ h=f\circ k$, entonces para $x\in S$ es
  % $f(h(x))=f(k(x))$ y por tanto $h(x)=k(x)$, por lo que $h=k$ y $f$ es un
  % monomorfismo. Veamos lo que ocurre en otros constructos.
  \begin{enumerate}
  \item En $\bSet$, por construcción, todo monomorfismo es inyectivo sin más que
    tomar $c=\{*\}$ en la definición anterior. Lo mismo ocurre en $\bTop$,
    $\bMet$, $\bOrd$ y otros constructos en que los elementos de un conjunto se
    identifiquen con los morfismos $\{*\}\to S$.
  \item En todos los constructos, los morfismos inyectivos son monomorfismos,
    pues si $f:X\to Y$ es un morfismo inyectivo y $h,k:S\to X$ con
    $f\circ h=f\circ k$, entonces para $x\in S$, $f(h(x))=f(k(x))$ y por tanto
    $h(x)=k(x)$, de modo que $h=k$.
  \item En $R\dash\bMod$, si $f:M\to N$ es un monomorfismo y $x,y\in M$ cumplen
    $f(x)=f(y)$, los morfismos $R\to M$ dados por $a\mapsto ax$ y $a\mapsto ay$
    quedan iguales al componerlos con $f$ y por tanto son iguales, por lo que
    $x=1x=1y=y$ y los monomorfismos son inyectivos.
  \item Aunque en la mayoría de constructos relevantes los monomorfismos son
    precisamente los morfismos inyectivos, esto no es siempre así, como muestra
    el constructo con objetos $\{*\}$, $\{\$\}$ y $\{x,y\}$ y como morfismos las
    identidades, las funciones con codominio $\{\$\}$ y el morfismo
    $*\mapsto x$, en el cual el morfismo $\{x,y\}\to\{\$\}$ es un monomorfismo
    pero no es inyectivo.
  \end{enumerate}
\end{example}

Hasta ahora para probar que los monomorfismos de un constructo son inyectivos
hemos tomado un objeto $D$ y un objeto $*\in D$ tal que, para cualquier otro
objeto $A$ y $a\in A$, podemos definir un morfismo $f:D\to A$ tal que $f(*)=a$.
Para las variedades algebraicas, sin embargo, no es inmediato ver cuál es este
objeto.

\begin{definition}
  Dados un conjunto graduado $\Omega=\{s_1,\dots,s_k\}$ donde cada $s_i$ tiene
  aridad $n_i$ y un conjunto finito $E$ de identidades en $\Omega$, el
  \conc{objeto libre} de $(\Omega,E)\dash\bAlg$ sobre el conjunto $X$ es el
  objeto construido tomando el conjunto de los árboles con raíz cuyos nodos se
  etiquetan con un $s_i$ y tienen $n_i$ hijos, o con un $x\in X$ y tienen 0
  hijos, haciendo el conjunto cociente por reescritura según las identidades en
  $E$, y definiendo las operaciones por construcción de árboles.
\end{definition}

\begin{example}[Objetos libres en variedades algebraicas comunes]\;
  \begin{enumerate}
  \item El monoide libre sobre $X$ está formado por las cadenas de elementos de
    $X$ junto a la concatenación, mientras que el semigrupo libre es similar
    pero excluyendo la cadena vacía.
  \item El grupo libre sobre $X$ está formado por cadenas de símbolos de la
    forma $x$ o $\overline x$ para $x\in X$, con la condición de que no aparecen
    subsecuencias $x\overline x$ u $\overline xx$, y la operación es la
    concatenación eliminando sucesivamente las subcadenas de esta forma.
  \item El grupo abeliano libre sobre $X$ es $\mathbb{Z}^X$.
  \item El anillo conmutativo libre sobre $\{x_1,\dots,x_n\}$ es el anillo de
    polinomios $\mathbb{Z}[x_1,\dots,x_n]$.
  \end{enumerate}
\end{example}

\begin{proposition}
  En las variedades algebraicas, los monomorfismos son precisamente los
  morfismos inyectivos.
\end{proposition}
\begin{proof}
  Sean $f:A\to B$ un monomorfismo en $(\Omega,E)\dash\bAlg$ y $x,y\in A$ con
  $f(x)=f(y)$. Si $F$ es la $(\Omega,E)$-álgebra libre sobre $\{*\}$, $h$ es el
  único morfismo $F\to A$ con $h(*)=x$ y $k:F\to A$ es el único con $h(*)=y$,
  para $c\in F$, $f(h(c))$ y $f(k(c))$ vienen dados por $f(h(*))$ y $f(k(*))$
  siguiendo la estructura de $c$, por lo que $f\circ h=f\circ k$, $h=k$ y,
  finalmente, $x=y$.
\end{proof}

\begin{proposition}\label{prop:mono-comp}\;
  \begin{enumerate}
  \item La composición de monomorfismos es un monomorfismo.
    \begin{proof}
      Sean $f:A\to B$ y $g:B\to C$ monomorfismos y $h,k:D\to A$ con
      $g\circ f\circ h=g\circ f\circ k$, entonces $f\circ h=f\circ k$ y por
      tanto $h=k$.
    \end{proof}
  \item Si $g\circ f$ es un monomorfismo, $f$ también lo es.
    \begin{proof}
      Si $g\circ f$ es un monomorfismo y $f\circ h=f\circ k$, entonces
      $g\circ f\circ h=g\circ f\circ k$ y por tanto $h=k$.
    \end{proof}
  \end{enumerate}
\end{proposition}

Para caracterizar los epimorfismos vemos que, en $\bSet$, si $f:A\to B$ no
es suprayectiva, sea $y\in B\setminus\Img{f}$, es posible crear dos funciones
$g,h:B\to\{0,1\}$ con $g\circ f=h\circ f$ pero $g(y)=0$ y $h(y)=1$, mientras
que fuera suprayectiva, $g\circ f=h\circ f$ implica claramente $g=h$.

\begin{definition}
  Un morfismo $f:a\to b$ es un \conc{epimorfismo}, $f:a\epicTo b$, si es
  cancelable por la derecha, es decir, si para $g,h:b\to c$ con
  $g\circ f=h\circ f$ se tiene $g=h$.
\end{definition}

La siguiente proposición se demuestra de forma similar que la
correspondiente para epimorfismos.

\begin{proposition}\label{prop:epi-comp}\;
  \begin{enumerate}
  \item La composición de epimorfismos es un epimorfismo.
  \item Si $g\circ f$ es un epimorfismo, $g$ también lo es.
  \end{enumerate}
\end{proposition}

\begin{example}\;
  \begin{enumerate}
  \item En los constructos, las funciones suprayectivas son epimorfismos, por el
    mismo argumento que en $\bSet$.
  \item En $\bSet$ los epimorfismos son precisamente las funciones
    suprayectivas, pues para $f:A\to B$ no suprayectiva es fácil encontrar
    $g,h:B\to\{0,1\}$ distintas con $g\circ f=h\circ f$. Lo mismo ocurre en
    $\bTop$ dotando a $\{0,1\}$ de la topología indiscreta y en $\bGrph$ usando
    el grafo completo de dos elementos.
  \item En $R\dash\bMod$, si $f:M\to N$ es un epimorfismo, sean
    $p,z:N\to\frac{N}{\Img{f}}$ la proyección canónica y la función constante en
    0, respectivamente, entonces $p\circ f=z\circ f$, luego $p=z$,
    $\frac{N}{\Img{f}}\cong0$ y así $\Img{f}=N$, con lo que los epimorfismos son
    suprayectivos.
  \item No en todas las variedades algebraicas los epimorfismos son son
    suprayectivos. Por ejemplo, en $\bRing$, la inclusión $u:\sInt\to\sRat$ es
    suprayectiva, pues si $f,g:\sRat\to R$ cumplen que
    $f\circ u=g\circ u:\sInt\to R$, entonces $f=g$ ya que para $x,y\in\sInt$,
    $f(\frac xy)=\frac{f(x)}{f(y)}=\frac{(f\circ u)(x)}{(f\circ u)(y)}$, y lo
    mismo ocurre con $g$.
  \item En $\bMet$ y $\bMetc$, los epimorfismos son precisamente las funciones
    con imagen densa.
    \begin{proof}
      Si $f:X\to Y$ es un morfismo con imagen densa y $g,h:Y\to Z$ son morfismos
      con $g\circ f=h\circ f$, para $y\in Y$ existe una sucesión
      $\{x_n\}_n\subseteq X$ con $y=\lim_nf(x_n)$, y por unicidad del límite y
      continuidad, $g(y)=\lim_ng(f(x_n))=\lim_nh(f(x_n))=h(y)$.  Recíprocamente,
      si $f:X\to Y$ no tiene imagen densa, existen $y_0\in Y$ y $r>0$ con
      $B(y_0,r)\cap\Img{f}=\emptyset$, por lo que $g,h:Y\mapsto\sReal$ dadas por
      $g(y)\coloneqq r$ y $h(y)\coloneqq\min\{d(y_0,y),r\}$ son retracciones
      continuas distintas con $g\circ f=h\circ f$.
    \end{proof}
  \end{enumerate}
\end{example}

Los dos últimos ejemplos muestran que, al contrario de lo que ocurre en muchos
campos del álgebra, no siempre los morfismos que son a la vez monomorfismos y
epimorfismos son isomorfismos.

\begin{definition}
  Un \conc{bimorfismo} es un morfismo que es a la vez un monomorfismo y un
  epimorfismo. Una categoría es \conc{equilibrada} si todo bimorfismo es un
  isomorfismo.
\end{definition}

\begin{example}\;
  \begin{enumerate}
  \item En una categoría fina, todos los morfismos son bimorfismos, pero en
    general no son isomorfismos.
  \item En un monoide visto como categoría, los bimorfismos son los elementos
    cancelables por ambos lados.
  \item Las identidades son bimorfismos, y en particular los isomorfismos son
    bimorfismos sin más que aplicar las proposiciones \ref{prop:mono-comp} y
    \ref{prop:epi-comp} a la composición de un isomorfismo con su inverso por
    ambos lados.
  \end{enumerate}
\end{example}

\section{Secciones y retracciones}

Hemos estudiado los morfismos cancelables por uno de los lados, por lo que cabe
preguntarse qué ocurre con los morfismos invertibles por uno de los
lados. Claramente este es un concepto más fuerte, y de hecho es estrictamente
más fuerte.

\begin{definition}
  Una \conc{sección} es un morfismo con inverso por la izquierda, y una
  \conc{retracción} es un morfismo con inverso por la derecha.  Es decir, si
  $f:a\to b$ y $g:b\to a$ son tales que $g\circ f=1_a$, entonces $f$ es una
  sección y $g$ es una retracción.
\end{definition}

Claramente las secciones son monomorfismos y las retracciones son
epimorfismos. Veamos algunos ejemplos.

\begin{example}\;
  \begin{enumerate}
  \item En $\bSet$, las secciones son las funciones inyectivas salvo las que van
    de un conjunto vacío a uno no vacío, y las retracciones son las funciones
    suprayectivas. Es fácil ver que esto último es equivalente al axioma de
    elección.
  \item En $R\dash\bMod$, las secciones son los monomorfismos cuya imagen es un
    sumando directo del codominio, y las retracciones son los epimorfismos cuyo
    núcleo es un sumando directo del dominio.
    \begin{proof}
      Si $f:M\to N$ es un monomorfismo cuya imagen es un sumando directo de $N$,
      por ejemplo $N=\Img{f}\oplus P$, un $g:N\to M$ que lleva los elementos de
      $\Img{f}$ a su única preimagen por $f$ y los $p\in P$ a 0 es inverso de
      $f$ por la derecha. Si es un epimorfismo cuyo núcleo es un sumando directo
      de $M$, por ejemplo $M=\ker{f}\oplus L$, sean $p:M\to\frac{M}{\ker{f}}$ la
      proyección canónica, $h:\frac{M}{\ker{f}}\cong L$ el isomorfismo <<obvio>>
      y $u:L\inTo M$ la inclusión, entonces $p\circ u\circ h=1$, pero el teorema
      del factor nos da un único isomorfismo
      $\overline{f}:\frac{M}{\ker{f}}\cong N$ con $\overline{f}\circ p=f$, por
      lo que $\overline{f}=f\circ u\circ h$ y
      $f\circ(u\circ h\circ\overline{f}^{-1})=1$.
      
      Para el recíproco, sean $f:M\monicTo N$ y $g:N\epicTo M$ con
      $g\circ f=1_M$, basta ver que $N=\Img{f}\oplus\ker{g}$. En efecto, todo
      $n\in N$ se descompone como suma de $n_1\coloneqq f(g(n))\in\Img{f}$ y
      $n_2\coloneqq n-f(g(n))\in\ker{g}$, y si $p_1\in\Img{f}$ y $p_2\in\ker{g}$
      cumplen $n=p_1+p_2$, sea $m_1\in M$ la preimagen de $p_1$ por $f$,
      entonces $g(n)=g(p_1+p_2)=g(f(m_1))+g(p_2)=m_1+0$ y
      $n_1=f(g(n))=g(m_1)=p_1$, con lo que $n_2=p_2$.
    \end{proof}
  \item En $K\dash\bVec$, como caso especial del apartando anterior, todos los
    monomorfismos son secciones y todos los epimorfismos son retracciones.
  \end{enumerate}
\end{example}

Las secciones y retracciones se conservan por composición del mismo modo que
lo hacen los monomorfismos y epimorfismos:

\begin{proposition}\;
  \begin{enumerate}
  \item La composición de secciones es una sección.
    \begin{proof}
      Si $f$ y $g$ son secciones con inversas por la izquierda respectivas
      $\overline f$ y $\overline g$, entonces
      $(\overline f\circ\overline g)\circ g\circ f=\overline f\circ f=1$.
    \end{proof}
  \item Si $g\circ f$ es una sección, $f$ es una sección.
    \begin{proof}
      Si $g\circ f$ es una sección con inversa por la izquierda $h$, entonces
      $h\circ g\circ f=1$, y $h\circ g$ es inversa por la izquierda de $f$.
    \end{proof}
  \end{enumerate}
\end{proposition}

De forma análoga se prueba la siguiente proposición.

\begin{proposition}\;
  \begin{enumerate}
  \item La composición de retracciones es una retracción.
  \item Si $g\circ f$ es una retracción, $g$ es una retracción.
  \end{enumerate}
\end{proposition}

Cabe preguntarse si un morfismo que es a la vez sección y retracción es
invertible. La respuesta es que sí, y además esta condición se puede relajar.

\begin{proposition}
  Para un morfismo $f$, las siguientes afirmaciones son equivalentes:
  \begin{enumerate}
  \item \label{enu:sr-iso} $f$ es un isomorfismo.
  \item \label{enu:sr-both} $f$ es una sección y una retracción.
  \item \label{enu:sr-relax-left} $f$ es un monomorfismo y una retracción.
  \item \label{enu:sr-relax-right} $f$ es una sección y un epimorfismo.
  \end{enumerate}
\end{proposition}
\begin{proof}
  Las implicaciones
  (\ref{enu:sr-iso})$\implies$(\ref{enu:sr-both})$\implies$(\ref{enu:sr-relax-left}),
  (\ref{enu:sr-relax-right}) son obvias.

  Para (\ref{enu:sr-both})$\implies$(\ref{enu:sr-iso}), si $f:a\to b$ y
  $g,h:b\to a$ son tales que $g\circ f=1_a$ y $f\circ h=1_b$, entonces
  $g=g\circ 1_b=g\circ f\circ h=1_a\circ h=h$, y $g=h$ es la inversa de $f$.

  Para (\ref{enu:sr-relax-left})$\implies$(\ref{enu:sr-iso}), sea $g$ un
  morfismo con $f\circ g=1$, entonces $f\circ g\circ f=1\circ f=f\circ 1$, y
  como $f$ es un monomorfismo, cancelando, $g\circ f=1$ y $f$ es un
  isomorfismo.

  De forma análoga (\ref{enu:sr-relax-right})$\implies$(\ref{enu:sr-iso}).
\end{proof}

\section{Dualidad}

La mayoría de las propiedades que hemos visto hasta ahora vienen en pares con
definiciones muy parecidas, y de hecho, cada vez que demostrábamos algo sobre
una de los propiedades, la misma idea servía para demostrar algo similar sobre
la otra. Esto ocurre mucho en teoría de categorías, y para sacar ventaja de esto
definimos el concepto de dualidad.

\begin{definition}
  Dada una categoría $\cC$, su \conc{categoría dual} es una categoría
  $\dual{\cC}$ con los mismos objetos y morfismos que $\cC$ y la misma función
  identidad pero tal que, para todo morfismo $f$,
  $\dom_{\dual{\cC}}{f}=\cod_{\cC}{f}$ y $\cod_{\dual{\cC}}{f}=\dom_{\cC}{f}$, y
  la composición se define como $f\circ_{\dual{\cC}}g\coloneqq g\circ_{\cC}f$.
\end{definition}

\begin{example}
  \begin{enumerate}
  \item Si $(X,\preceq)$ es un conjunto preordenado, su dual visto como
    categoría es $(X,\succeq)$.
  \item El dual de una categoría discreta es ella misma.
  \end{enumerate}
\end{example}

La principal utilidad de la categoría dual está en definir el dual de un
concepto o propiedad.

\begin{definition}
  Si $P$ es un predicado aplicable a una o más categorías $\cC_1,\dots,\cC_n$,
  su \conc{dual} es
  $\dual{P}(\cC_1,\dots,\cC_n)\equiv P(\dual{\cC_1},\dots,\dual{\cC_n})$, y
  decimos que $P$ es \conc{auto-dual} si $P\equiv\dual{P}$.
\end{definition}

En esencia, tomar la categoría dual consiste en invertir el sentido de las
flechas en la categoría, y tomar el predicado dual consiste en invertir el
sentido de los morfismos que se mencionan. En general al tratar un concepto
categórico es conveniente fijarse en el dual del concepto, ya que suele ser un
concepto relevante y, además, las propiedades de dicho concepto se derivan
directamente de las del concepto original, sin necesidad de demostración aparte.

\section{Producto y coproducto}

En muchas categorías es posible tomar el producto de dos o más objetos, como el
de dos conjuntos, el espacio topológico producto, el espacio vectorial producto,
etc. Para definir productos en una categoría arbitraria tenemos que encontrar
una propiedad universal del producto, de modo que el producto quede definido de
forma única salvo homomorfismo.

Para ello, si, por ejemplo, $A$ y $B$ son conjuntos, el producto $A\times B$ es
un conjunto de pares de un elemento $a\in A$ y un elemento $b\in B$, de modo que
dadas dos funciones $f:X\to A$ y $g:X\to B$ existe una única función
$h:X\to (A\times B)$ tal que, para todo $x\in X$, $h(x)=(f(x),g(x))$. Dicho de
otro modo, si $p:A\times B\to A$ y $q:A\times B\to B$ son las proyecciones sobre
las componentes, $h$ es la única función con $f=p\circ h$ y $g=q\circ h$. Esta
definición se puede generalizar a categorías arbitrarias y a una cantidad
arbitraria de factores, y la figura \ref{fig:product} muestra el caso para dos
factores.

\begin{figure}
  \centering
  \begin{diagram}
    \path                    (2,2) node(X){$x$}
          (0,0) node(A){$a$} (2,0) node(AB){$a\times b$} (4,0) node(B){$b$};
    \draw[->] (AB) -- node[below]{$p$} (A);
    \draw[->] (AB) -- node[below]{$q$} (B);
    \draw[->] (X) -- node[left]{$f$} (A);
    \draw[->] (X) -- node[right]{$g$} (B);
    \draw[->,dotted] (X) -- node{$f\times g$} (AB);
  \end{diagram}
  \caption[Objeto producto]{Objeto producto de $a$ y $b$. Para
    cualesquiera $x$, $f$ y $g$, existe un único $f\times g$ de modo
    que el diagrama conmuta.}
  \label{fig:product}
\end{figure}

% Para tratar productos es conveniente ver primero el concepto de fuente.

% \begin{definition}
%   Una \conc{fuente} es un par $(b,(f_i:b\to a_i)_{i\in I})$ formada por un
%   objeto $b$ y una familia de morfismos con dominio $b$. Llamamos \conc{dominio}
%   de la fuente a $b$ y \conc{codominio} a $(a_i)_{i\in I}$, y normalmente
%   nos referimos a la fuente por su familia de morfismos $(f_i)_{i\in I}$.
% \end{definition}

\begin{definition}
  Un objeto $b$ es un \conc{producto} de una familia de objetos $(a_i)_{i\in I}$
  si existe una familia de morfismos $(p_i:b\to a_i)_{i\in I}$, llamados
  \conc{proyecciones}, tales que para cualquier familia de morfismos
  $(f_i:x\to a_i)_{i\in I}$ existe un único $f:x\to b$ tal que $f_i=p_i\circ f$
  para cada $i$. Llamamos producto a $b$ y a la familia de proyecciones
  indistintamente.
\end{definition}

\begin{example}\;
  \begin{enumerate}
  \item En $\bSet$, toda familia pequeña de conjuntos tiene un producto, el
    producto directo habitual. Lo mismo ocurre en $R\dash\bMod$ y en particular
    en $\bVec$ y en $\bAb$, así como en $\bGrp$ con el producto de grupos y en
    $\bTop$ con el producto topológico.
  \item Todo objeto es el producto de sí mismo tomando como proyección el
    morfismo identidad. De hecho, un objeto $a$ es el producto de la familia
    unipuntual $(b)$ si y sólo si $a\cong b$.
  \item Un objeto es producto de la familia vacía si y sólo si es final.
  \item En un conjunto preordenado visto como categoría, el producto de una
    familia de elementos es su supremo, si existe.
  \end{enumerate}
\end{example}

\begin{proposition}
  El producto de una familia de objetos es único salvo isomorfismo.
\end{proposition}
\begin{proof}
  Sean $b$ y $c$ productos de $(a_i)_{i\in I}$ con familias de proyecciones
  $(f_i:b\to a_i)_{i\in I}$ y $(g_i:c\to a_i)_{i\in I}$, existe $h:b\to c$ tal
  que cada $f_i=g_i\circ h$ y $k:c\to b$ tal que cada $g_i=f_i\circ k$, pero
  entonces $g_i=g_i=g_i\circ h\circ k$ para cada $i$, y como para la familia
  $(g_i)_i$ debe haber un único $f:c\to c$ con cada $g_i=g_i\circ f$, debe ser
  $h\circ k=f=1_c$, y análogamente $k\circ h=1_b$. Es fácil ver que los
  isomorfismos conservan productos.
\end{proof}

En vista de esto, llamamos $\prod_{i\in I}a_i$ al objeto producto de
$(a_i)_{i\in I}$. Denotamos el objeto producto de dos objetos $a$ y $b$ como
$a\times b$, y esta notación se puede extender a un número finito de factores
($a_1\times\dots\times a_n$) debido a que el producto es asociativo, en el
sentido de la siguiente proposición.

\begin{proposition}
  Si $(p_i:c\to b_i)_{i\in I}$ es un producto y, para cada $i$,
  $(q_{ij}:b_i\to a_{ij})_{j\in J_i}$ es un producto, entonces
  $(q_{ij}\circ p_i:c\to a_{ij})_{i\in I}^{j\in J_i}$ es un producto.
\end{proposition}
\begin{proof}
  Sea $(f_{ij}:x\to a_{ij})_{ij}$ una familia de morfismos, para cada $i\in I$
  existe $g_i:x\to b_i$ tal que $f_{ij}=q_{ij}\circ g_i$ para todo $j$, por lo
  que existe $h:x\to c$ tal que $g_i=p_i\circ h$ para todo $i$ y este cumple
  $f_{ij}=q_{ij}\circ p_i\circ h$ para todo $i,j$. Además, los $g_i$ y $h$ son
  únicos, por lo que si hubiera otro $k:x\to c$ con
  $f_{ij}=q_{ij}\circ p_i\circ k$ para todo $i,j$, necesariamente cada
  $p_i\circ k=g_i$ y por tanto $k=h$.
\end{proof}

Así, en particular, $(a\times b)\times c=a\times b\times c$ y una categoría
tiene todos los productos finitos si y sólo si tiene objetos terminales y
productos de todos los pares de objetos.

\begin{definition}
  Dado un objeto $a$ y un conjunto $I$, llamamos \conc{$I$-ésima potencia de
    $a$} a $a^I\coloneqq\prod_{i\in I}a$.
\end{definition}

% Puntos de expansión: prop. 10.38 (caracterización de coseparadores según
% potencias), 10.28 (las proyecciones suelen ser retracciones), 10.34 (productos
% de morfismos arbitrarios), mono-fuentes y mono-fuentes extremas.

El concepto dual del producto es el coproducto, que se muestra gráficamente en
la figura \ref{fig:coproduct}.

\begin{figure}
  \centering
  \begin{diagram}
    \path (0,2) node(A){$a$} (2,2) node(AB){$a\times b$} (4,2) node(B){$b$}
                             (2,0) node(X){$x$};
    \draw[->] (A) -- node[above]{$u$} (AB);
    \draw[->] (B) -- node[above]{$v$} (AB);
    \draw[->] (A) -- node[left]{$f$} (X);
    \draw[->] (B) -- node[right]{$g$} (X);
    \draw[->,dotted] (AB) -- node{$f\oplus g$} (X);
  \end{diagram}
  \caption[Objeto coproducto]{Objeto coproducto de $a$ y $b$. Para cualesquiera
    $X$, $f$ y $g$, existe un único $f\oplus g$ de modo que el
    diagrama conmuta.}
  \label{fig:coproduct}
\end{figure}

\begin{definition}
  Un objeto $c$ es un \conc{coproducto} de una familia de objetos
  $(a_i)_{i\in I}$ si existe una familia de morfismos $(u_i:a_i\to c)_{i\in I}$,
  llamados \conc{inyecciones}, tales que para cualquier familia de morfismos
  $(f_i:a_i\to x)_{i\in I}$ existe un único $f:c\to x$ tal que $f_i=f\circ u_i$
  para cada $i$. Llamamos coproducto tanto a $c$ como a la familia de
  inyecciones.
\end{definition}

\begin{example}
  \begin{enumerate}
  \item En $\bSet$, el coproducto de una familia pequeña de conjuntos (pequeños)
    $(a_i)_{i\in I}$ es la \conc{unión disjunta},
    $\biguplus_{i\in I}a_i\coloneqq\bigcup_{i\in I}(a_i\times\{i\})$, aunque si
    los $a_i$ son disjuntos dos a dos se puede tomar como coproducto la unión
    convencional.
  \item En $\bTop$, el coproducto de una familia pequeña de espacios topológicos
    es el espacio topológico de la unión disjunta, que tiene como abiertos las
    uniones de abiertos de cada espacio.
  \item En $R\dash\bMod$, el coproducto de una familia pequeña de módulos
    $\{a_i\}_{i\in I}$ es la suma directa $\bigoplus_{i\in I}a_i$, el subespacio
    del producto directo $\prod_{i\in I}a_i$ formado por los elementos con casi
    todas las entradas nulas. En particular, si $I$ es finito el producto
    coincide con el coproducto.
  \item Todo objeto es coproducto de sí mismo.
  \item Un objeto es coproducto de la familia vacía si y sólo si es final.
  \item En un conjunto preordenado visto como categoría, el coproducto de una
    familia de elementos es su ínfimo, si existe.
  \end{enumerate}  
\end{example}

Las siguientes propiedades son duales de las correspondientes del producto.

\begin{proposition}
  El coproducto  de una familia de objetos es único salvo isomorfismo.
\end{proposition}

Esto nos permite llamar $\coprod_{i\in I}a_i$ al objeto coproducto de
$(a_i)_{i\in I}$, y $a\oplus b$ al objeto coproducto de dos objetos $a$
y $b$.

\begin{proposition}
  Si $(u_i:b_i\to c)_{i\in I}$ es un coproducto y, para cada $i$,
  $(v_{ij}:a_{ij}\to b_i)_{j\in J_i}$ es un coproducto, entonces
  $(u_i\circ v_{ij}:a_{ij}\to c)_{i\in I}^{j\in J_i}$ es un coproducto.
\end{proposition}

\begin{definition}
  Dado un objeto $a$ y un conjunto $I$, llamamos \conc{$I$-ésima copotencia de
    $a$} a $^Ia\coloneqq\coprod_{i\in I}a$.
\end{definition}

\section{Núcleo y conúcleo}

En álgebra es común hablar del núcleo de un morfismo $f:a\to b$ como el conjunto
de puntos $x\in a$ con $f(x)=0$, que es un subobjeto de $a$.  Entonces, si
$k\subseteq a$ es el núcleo de $f$ y $u:k\inTo a$ es la inclusión, se tiene
$f\circ u=0$ y, si $u':k'\to a$ es otro morfismo tal que $f\circ u'=0$, existe
$\tilde u:k'\to k$ tal que $u'=u\circ\tilde u$. Esta definición nos sirve para
teoría de categorías salvo por la presencia del 0.

Si, para cada objeto $a$, llamamos $i_a:0\to a$ a la única flecha desde el
objeto cero hasta $a$ y llamamos $t_a:a\to 0$ a la única flecha desde $a$ hasta
el objeto 0, entonces podríamos escribir <<$f\circ u=0$>> como
<<$f\circ u=i_b\circ t_k$>>, pero esta caracterización sólo nos sirve para el
caso de categorías con cero. En general, podemos notar que $t_k=t_a\circ u$, de
modo que $u$ <<iguala>> $f$ con $i_b\circ t_a$, y aunque no podemos hablar del
núcleo de una función, podemos hablar del núcleo de dos funciones en el sentido
siguiente.

\begin{definition}
  Dados dos morfismos $f,g:a\to b$, un morfismo $e:k\to a$ es un \conc{núcleo}
  de $f$ y $g$ si y sólo si $f\circ e=g\circ e$ y, para todo $e':k'\to a$ con
  $f\circ e'=g\circ e'$, existe un único $\tilde e:k'\to k$ con
  $e'=e\circ\tilde e$, es decir, tal que la figura \ref{fig:equalizer} conmuta.
  Si existe un cero $0$, un \conc{núcleo} de $f:a\to b$ es un núcleo de $f$ y la
  composición $a\to 0\to b$.
\end{definition}

\begin{figure}
  \centering
  \begin{diagram}
    \path (0,2) node(KP){$k'$}
                (1.2,1.2) node{$e'$}
          (0,0) node(K){$k$} (2,0) node(A){$a$} (4,0) node(B){$b$};
    \draw[->,dotted] (KP) -- node[left]{$\tilde e$} (K);
    \draw[->] (KP) -- (A);
    \draw[->] (K) -- node[below]{$e$} (A);
    \draw[->] (A.15) -- node[above]{$f$} (B.165);
    \draw[->] (A.345) -- node[below]{$g$} (B.195);
  \end{diagram}
  \caption[Núcleo de dos homomorfismos]{Núcleo $e$ de un par de homomorfismos $f$ y $g$.}
  \label{fig:equalizer}
\end{figure}

\begin{example}\;
  \label{ex:equalizer}
  \begin{enumerate}
  \item En $\bSet$, para $f,g:a\to b$, $k\coloneqq\{x\in a\mid f(x)=g(x)\}$ es
    un subconjunto de $a$ y la inclusión $u:k\inTo a$ es un núcleo de $f$ y
    $g$. Lo mismo ocurre en $R\dash\bMod$, $\bTop$, $\bGrp$ y $\bGrph$ dotando a
    $k$ de su estructura como submódulo, subespacio topológico, subgrupo o
    subgrafo.
  \item En $\bRng$, al contrario que en muchas estructuras algebraicas, el
    concepto de núcleo de un morfismo en teoría de categorías que en álgebra no
    coincide con el convencional, ya que convencionalmente el núcleo de un
    anillo no es un anillo sino un ideal.
  \item En $R\dash\bMat$, dadas dos matrices $n\times m$ $A,B:m\to n$, un núcleo
    de $A$ y $B$ es precisamente una matriz $X:p\to m$ cuyas columnas forman una
    base de soluciones de la ecuación $Ax=Bx$. En efecto, $AX=BX$ y, si
    $Y:q\to m$ cumple que $AY=BY$, cada columna de $Y$ es combinación lineal
    única de las columnas de $X$ y hay por tanto existe una única $Z$ con
    $Y=XZ$. La existencia de $Z$ implica que las columnas de $X$ deben generar
    el espacio de soluciones, y la unicidad implica que estan deben ser
    linealmente independientes. % Ref. by siguiente párrafo ("último ejemplo")
  \end{enumerate}
  % TODO Ejemplo en (Ω,E)-Alg
\end{example}

En el último ejemplo hemos demostrado que todos los núcleos tienen la misma
forma. Esto no es necesario, y basta con encontrar un núcleo, por la siguiente
definición.

\begin{proposition}
  \label{prop:equ-uniq}
  El núcleo es esencialmente único, es decir, si $e:k\to a$ es un núcleo de
  $f,g:a\to b$, el resto de núcleos de $f$ y $g$ son precisamente los morfismos
  de la forma $e\circ h:k'\to a$ donde $h:k'\to k$ es un isomorfismo.
\end{proposition}
\begin{proof}
  Claramente $e\circ h$ es un núcleo, pues $f\circ e\circ h=g\circ e\circ h$ y,
  si $e'$ cumple $f\circ e'=g\circ e'$ y $\tilde e$ es el único morfismo con
  $e'=e\circ\tilde e$, entonces $h^{-1}\circ\tilde e$ es el único con
  $e'=(e\circ h)\circ(h^{-1}\circ\tilde e)$. Para el recíproco, sean $e:k\to a$
  y $e':k'\to a$ núcleos de $f$ y $g$, existen un único $h:k'\to k$ con
  $e'=e\circ h$ y un único $h':k\to k'$ con $e=e'\circ h'$, de modo que
  $e\circ 1_k=e\circ h\circ h'$ y, por unicidad, $1_k=h\circ h'$, pero
  análogamente $1_{k'}=h'\circ h$, luego $h$ es un isomorfismo.
\end{proof}

Esta prueba parece sugerir que los núcleos son monomorfismos. Parece interesante
entonces comparar el concepto de núcleo con el de monomorfismo, así como con el
concepto relacionado de sección.

\begin{proposition}\;
  \label{prop:equ-middle}
  \begin{enumerate}
  \item Todo núcleo es un monomorfismo.
    \begin{proof}
      Si $e:k\to a$ es el núcleo de $f,g:a\to b$ y $h,k:c\to k$ cumplen
      $e\circ h=e\circ k\eqqcolon e'$, entonces $f\circ e'=g\circ e'$, luego el
      $h$ con $e'=e\circ h$ es único y por tanto $h=k$.
    \end{proof}
  \item Toda sección es un núcleo. En concreto, si $f:a\to b$ es una sección y
    $g:b\to a$ es la correspondiente retracción, entonces $f$ es núcleo de
    $f\circ g$ y $1_b$.
    \begin{proof}
      Para empezar, $(f\circ g)\circ f=f\circ(g\circ f)=f=1_b\circ f$, pero si
      $e:k\to b$ es tal que $(f\circ g)\circ e=1_b\circ e$, entonces $g\circ e$
      es el morfismo con $f\circ(g\circ e)=e$, y es único porque $f$ es un
      monomorfismo.
    \end{proof}
  \item \label{enu:equ-mid-strict} Los recíprocos no se cumplen.
    \begin{proof}
      En $\bAb$, el único homomorfismo $\sInt\to\sInt_2$ tiene como núcleo la
      inclusión $2\sInt\inTo\sInt$, que claramente no es una seccion (no tiene
      inverso por la izquierda). Y en $\bRng$, la inclusión $\sInt\inTo\sRat$ es
      un monomorfismo pero no es un núcleo, ya que de serlo, como también es un
      epimorfismo, sería un isomorfismo por la siguiente proposición
      (\ref{prop:iso-equalizer}).
      % TODO Igual debería poner la "siguiente" proposición antes.
    \end{proof}
  \end{enumerate}
\end{proposition}

\begin{proposition}
  \label{prop:iso-equalizer} Si $e:k\to a$ es núcleo de $f,g:a\to b$, son
  equivalentes:
  \begin{enumerate}
  \item \label{enu:eq-equal} $f=g$.
  \item \label{enu:eq-ident} $1_a$ es núcleo de $f$ y $g$.
  \item \label{enu:eq-iso} $e$ es un isomorfismo.
  \item \label{enu:eq-epi} $e$ es un epimorfismo.
  \end{enumerate}
\end{proposition}
\begin{proof}
  Claramente (\ref{enu:eq-equal})$\implies$(\ref{enu:eq-ident});
  (\ref{enu:eq-ident})$\implies$(\ref{enu:eq-iso}) por (\ref{prop:equ-uniq});
  obviamente (\ref{enu:eq-iso})$\implies$(\ref{enu:eq-epi}), y
  (\ref{enu:eq-epi})$\implies$(\ref{enu:eq-equal}) por definición de
  núcleo y de epimorfismo.
\end{proof}

El concepto dual de núcleo es el de conúcleo.

\begin{definition}
  Dados dos morfismos $f,g:a\to b$, un morfismo $c:b\to q$ es un \conc{conúcleo}
  de $f$ y $g$ si y sólo si $c\circ f=c\circ g$ y, para todo $c':b\to q'$ con
  $c'\circ f=c'\circ g$, existe un único $\overline c:q\to q'$ con
  $c'=\overline c\circ c$, es decir, tal que la figura \ref{fig:coequalizer}
  conmuta. Si existe un cero 0, un \conc{conúcleo} de $f:a\to b$ es un conúcleo
  de $f$ y la composición $a\to 0\to b$.
\end{definition}

\begin{figure}
  \centering
  \begin{diagram}
    \path (0,-2) node(KP){$q'$}
                  (-1.2,-1.2) node{$c'$}
          (0,0) node(K){$q$} (-2,0) node(A){$b$} (-4,0) node(B){$a$};
    \draw[<-,dotted] (KP) -- node[right]{$\overline c$} (K);
    \draw[<-] (KP) -- (A);
    \draw[<-] (K) -- node[above]{$c$} (A);
    \draw[<-] (A.165) -- node[above]{$f$} (B.15);
    \draw[<-] (A.195) -- node[below]{$g$} (B.345);
  \end{diagram}
  \caption[Conúcleo de dos homomorfismos]{Conúcleo $q$ de un par de homomorfismos $f$ y $g$.}
  \label{fig:coequalizer}
\end{figure}

\begin{example}\;
  \begin{enumerate}
  \item Sean $f,g:a\to b$ en $\bSet$ y sea $\sim$ la menor relación de
    equivalencia con $f(x)\sim g(x)$ para todo $x\in a$. Entonces la proyección
    al conjunto cociente $c:b\to\frac{b}{\sim}$ es un conúcleo de $f$ y $g$.
    \begin{proof}
      Claramente $c\circ f=c\circ g$ y, si $c':b\to r$ es tal que
      $c'\circ f=c'\circ g$, entonces $\overline c:\frac{b}{\sim}\to r$ dada por
      $\overline c(\overline x)=c'(x)$ es la única función con
      $\overline c\circ c=c'$, y queda ver que está bien definida. Pero si
      $x\sim{}y$, bien $x=y$, bien hay una cadena $x=t_0,\dots,t_k=y\in b$ y
      $s_1,\dots,s_k\in a$ con $(t_{i-1},t_i)$ igual a $(f(s_i),g(s_i))$ o a
      $(g(s_i),f(s_i))$ para cada $i$, pero entonces por hipótesis
      $c'(t_{i-1})=c'(t_i)$ para cada $i$ y por tanto $c'(x)=c'(y)$.
    \end{proof}
  \item En $\bTop$, los conúcleos se calculan como en $\bSet$ asignando a
    $\frac{b}{\sim}$ la topología cociente, cuyos abiertos son los subconjuntos
    de $\frac{b}{\sim}$ con $c^{-1}\subseteq b$ abierto. Es fácil ver que esta
    topología hace a $c$ y $\overline c$ del apartado anterior continuas
    (supuesto que $c'$ es continua).
  \item En $\bGrph$, ocurre lo mismo, estableciendo como ejes en
    $\frac{b}{\sim}$ las imágenes de ejes en $b$. En $\bPrord$ los conúcleos se
    construyen de igual forma pero tomando la clausura transitiva de la relación
    resultante en $\frac{b}{\sim}$.
  \item En $\bRng$, el conúcleo de $f,g:A\to B$ es la proyección canónica
    $p:B\to\frac{B}{I}$, donde $I\coloneqq(\Img{f-g})$ es el ideal generado por
    la imagen de $f-g$.
    \begin{proof}
      Para $a\in A$, $p(f(a))-p(g(a))=p((f-g)(a))=0$, pues $(f-g)(a)\in\ker p$,
      y si $c:B\to C$ es un morfismo que cumple $c\circ f=c\circ g$ y por tanto
      $c\circ (f-g)=0$, y los elementos de $I$ son combinaciones lineales de
      imágenes de $f-g$, se tiene $I\subseteq\ker c$ y, por el teorema del
      factor, existe un único $\overline c:\frac{B}{I}\to C$ con
      $c=\overline c\circ p$.
    \end{proof}
  \item En $(\Omega,E)\dash\bAlg$ el conúcleo se establece de la misma forma
    pero sustituyendo $\sim$ por la menor relación de congruencia en $b$ con
    $f(x)\sim g(x)$ para todo $x\in a$. Una \conc{relación de congruencia} es
    una relación de equivalencia en que, para cada operación $f:b^n\to b$ en la
    acción de $\Omega$ asociada a $b$ y cada $x_1,\dots,x_n,y_1,\dots,y_n\in b$
    se tiene que, si cada $x_i\sim y_i$, entonces
    $f(x_1,\dots,x_n)\sim f(y_1,\dots,y_n)$. Esta propiedad permite dotar a
    $\frac{b}{\sim}$ una estructura algebraica cuyas operaciones se derivan de
    las de $b$ de la forma evidente.
  \end{enumerate}
\end{example}

Las siguientes proposiciones son las duales de las vistas para el núcleo.

\begin{proposition}
  El conúcleo es esencialmente único, es decir, si $c:b\to q$ es un conúcleo de
  $f,g:a\to b$, el resto de conúcleos de $f$ y $g$ son precisamente los
  morfismos de la forma $h\circ c:b\to q'$ donde $h:q\to q'$ es un isomorfismo.
\end{proposition}

\begin{proposition}\;
  \begin{enumerate}
  \item Todo conúcleo es un epimorfismo.
  \item Toda retracción es un conúcleo. En concreto, si $g:a\to b$ es una
    retracción y $f:b\to a$ es la correspondiente sección, entonces $g$ es
    conúcleo de $f\circ g$ y $1_a$.
  \item Los recíprocos no se cumplen.
  \end{enumerate}
\end{proposition}

\begin{proposition}
  Si $c:b\to q$ es conúcleo de $f,g:a\to b$, son equivalentes:
  \begin{enumerate}
  \item $f=g$.
  \item $1_b$ es conúcleo de $f$ y $g$.
  \item $c$ es un isomorfismo.
  \item $c$ es un epimorfismo.
  \end{enumerate}
\end{proposition}

\section{Subobjetos y objetos cociente}

En los ejemplos de núcleos y conúcleos que hemos visto, por lo general los
núcleos son subobjetos del dominio de las funciones involucradas, mientras que
los conúcleos son objetos cociente. En teoría de categorías, como las categorías
no tienen por qué ser conjuntos, resulta difícil definir estos conceptos, y de
hecho no hay una sóla definición para todos los casos, aunque todas las
definiciones se basan en una estructura común.

\begin{definition}
  Sea $M$ un conjunto de monomorfismos. Un \conc{$M$-subobjeto} de un objeto $a$
  es un par $(b,m)$ formado por un objeto $b$ y un morfismo $m:b\to a$ en $M$.
\end{definition}

Las distintas definiciones dependen, pues, de qué conjunto $M$ usemos. Si
estamos en un constructo tiene sentido considerar el conjunto de morfismos que
son inclusiones. Para el caso general, sin embargo, existen varios conceptos
usados en la práctica, siendo el más amplio el de monomorfismo y el más
restringido el de sección. Por supuesto, uno de los conceptos intermedios, útil
para categorías algebraicas como $R\dash\bMod$, es el de un monomorfismo que es
el núcleo de algún morfismo.

\begin{definition}
  Un monomorfismo es \conc{regular} si es el núcleo de algún par de morfismos.
\end{definition}

Así, si $M$ es el conjunto de los monomorfismos regulares, los $M$-subobjetos se
llamarían \conc{subobjetos regulares}, mientras que si $M$ es el conjunto de
todos los monomorfismos, los $M$-subobjetos son los \conc{subobjetos} a secas.

Para todos estos conceptos aplicables en general, el conjunto de subobjetos es
cerrado para isomorfismos, en el sentido siguiente.

\begin{definition}
  Dados dos $M$-subobjetos $(a,m)$ y $(b,n)$ de un objeto $c$, $(a,m)$ es
  \conc{más pequeño} que $(b,n)$, $(a,m)\leq(b,n)$, si existe $f:a\to b$ con
  $m=n\circ f$, y $(a,m)$ y $(b,n)$ son \conc{isomorfos}, $(a,m)\cong(b,n)$, si
  además $f$ es un isomorfismo.
\end{definition}

Cabe destacar la importancia de los morfismos en la definición de
subobjetos. Por ejemplo, en $\bOrd$, $\{1,3,5\}$ y $\{2,4,6\}$ son subconjuntos
ordenados de $\sNat$ y, vistos como objetos, son isomorfos y por tanto
<<esencialmente lo mismo>>; lo que los distingue es el monomorfismo inclusión en
$\sNat$. A lo largo de la teoría de categorías se da mucho esta situación, en
que los objetos como tales, <<vistos desde fuera>>, no dan demasiada información
y, sin embargo, los morfismos permiten codificar la información relevante.

\begin{example}\;
  \label{ex:reg-mono}
  \begin{enumerate}
  \item \label{enu:reg-mono-set} En $\bSet$, los monomorfismos regulares son las
    funciones inyectivas, por lo que todos los monomorfismos son regulares. En
    efecto, las funciones inyectivas son aquellas isomorfas a una inclusión, y
    toda inclusión $u:E\inTo A$ es núcleo del par de funciones $A\to\{0,1\}$ en
    el que una lleva todos los elementos a 0 y otra lleva a 0 los elementos de
    $E$ y a 1 el resto. Por tanto los subobjetos regulares se identifican (salvo
    isomorfismo) con los subconjuntos, y un subobjeto es más pequeño que otro si
    y sólo si es un subconjunto del otro.
  \item En $\bTop$ los monomorfismos regulares son, salvo isomorfismo, las
    inclusiones de subespacios, sin más que tomar la prueba del apartado
    anterior y dotar a $\{0,1\}$ de la topología indiscreta, por lo que los
    monomorfismos regulares son los subespacios topológicos.
  \item En muchas categorías algebraicas como $R\dash\bMod$ o $\bGrp$, todos los
    monomorfismos son regulares, de modo que los subobjetos regulares son
    respectivamente los submódulos y los grupos. En el caso de $R\dash\bMod$, si
    $e:K\monicTo M$ es un monomorfismo de módulos, $e$ es el núcleo de la
    proyección canónica $M\mapsto\frac{M}{\Img{e}}$, y en $\bGrp$ la prueba es
    similar.
  \item En $R\dash\bMod$, todos los monomorfismos son regulares. En efecto, un
    monomorfismo $m:M\to N$ es núcleo de la proyección canónica
    $p:N\to\frac{N}{\Img{m}}$ y el correspondiente morfismo cero.
  \item En $\bRing$ hemos visto (\ref{prop:equ-middle}, apartado
    \ref{enu:equ-mid-strict}) que no todos los monomorfismos son regulares. Sin
    embargo, los subobjetos (a secas) son precisamente los subanillos.
  \item En una categoría fina, sólo las identidades son regulares.
  \item En $R\dash\bMat$, los núcleos son las matrices $m\times p$, $p\leq m$,
    de rango máximo (\ref{ex:equalizer}), con lo que los subobjetos de un número
    $m$ son los pares formados por un $p\leq m$ y una matriz $m\times p$ de
    rango $p$.
  \end{enumerate}
\end{example}

El concepto dual de subobjeto es el de objeto cociente.

\begin{definition}
  Sea $E$ un conjunto de epimorfismos. Un \conc{$E$-objeto cociente} de un
  objeto $a$ es un par $(b,e)$ formado por un objeto $b$ y un morfismo
  $e:a\to b$ en $E$.
\end{definition}

\begin{definition}
  Un epimorfismo es \conc{regular} si es el conúcleo de algún par de morfismos.
\end{definition}

Si $E$ es el conjunto de los epimorfismos regulares, hablamos de \conc{objetos
  cociente regulares}, mientras que si $E$ es el conjunto de todos los
epimorfismos hablamos de \conc{objetos cociente} a secas.

Si estamos en constructos tiene sentido considerar los epimorfismos que son
proyecciones al conjunto cociente por alguna relación de equivalencia en el
dominio, o a algún conjunto irredundante de representantes de dicha relación. En
la mayoría de constructos relevantes, si para una cierta relación de
equivalencia existe una de estas proyecciones, el resto también existe y son
isomorfas en el sentido siguiente.

\begin{definition}
  Dados dos $E$-objetos cociente $(b,d)$ y $(c,e)$ de un objeto $a$, $(b,d)$ es
  \conc{más grande} que $(c,e)$, $(b,d)\geq(c,e)$, si existe $f:b\to c$ con
  $e=f\circ d$, y $(b,d)$ y $(c,e)$ son \conc{isomorfos}, $(b,d)\cong(c,e)$, si
  además $f$ es un isomorfismo.
\end{definition}

\begin{example}\;
  \begin{enumerate}
  \item \label{enu:reg-epi-set} En $\bSet$, todos los epimorfismos son
    regulares, pues si $e:A\to B$ es suprayectiva, es el conúcleo de las dos
    proyecciones $D_e\to A$ con
    $D_e\coloneqq\{(a,a')\in A\times A\mid e(a)=e(a')\}$, por un argumento
    similar al usado en el ejemplo \ref{ex:reg-mono}, apartado
    \ref{enu:reg-mono-set}. Además, claramente los objetos cociente de $A$ son,
    salvo isomorfismo, los conjuntos cociente con sus proyecciones, y si $(B,d)$
    y $(C,e)$ son objetos cociente (regulares) de $A$, entonces $(B,d)\geq(C,e)$
    si y sólo si $D_d\subseteq D_e$.
  \item De forma similar se ve que los objetos cocientes regulares en $\bTop$
    son, salvo isomorfismo, los objetos topológicos cociente y las
    correspondientes proyecciones. Aquí, sin embargo, no todos los epimorfismos
    son regulares, pues el codominio debe tener la topología <<correcta>>, y no
    una más gruesa.
  \item En $R\dash\bMod$, todos los epimorfismos son regulares. En efecto, si
    $e:M\to N$ es un morfismo, $D_e$ (apartado \ref{enu:reg-epi-set}) es el
    núcleo del homomorfismo $(m,m')\mapsto e(m-m')$, por lo que es un $R$-módulo
    y, por el argumento del apartado \ref{enu:reg-epi-set}, $e$ es el conúcleo
    de las dos proyecciones $D_e\to M$. Además es fácil comprobar que los
    objetos cociente son los módulos cociente.
  \end{enumerate}
\end{example}

%%% Local Variables:
%%% mode: latex
%%% TeX-master: "main"
%%% End:
